\section{Introduction}
\label{sec:intro}

\subsection{The IDR framework and its limitations}

The Irreproducible Discovery Rate framework of \citet{LiBrownHuangBickel2011} has become a standard tool for replicate-experiment reproducibility quantification. It posits a two-component mixture in which the reproducible component is generated by a bivariate latent Gaussian model with shared mean and variance across replicates, fits the parameters by an EM procedure on pseudo-data obtained from a marginal probability-integral transform, and reports a per-feature local idr together with a Sun--Cai step-up cutoff that controls the marginal false discovery rate. The framework is rank-invariant, scale-invariant, and threshold-free, which makes it a natural choice for consortium pipelines that aggregate heterogeneous platforms.

The original framework, however, has three rigid assumptions that modern multi-replicate designs can violate. First, the construction is bivariate: exactly two replicates ($K = 2$) are accommodated, while modern ENCODE 4 transcription-factor experiments typically use $K \in \{2, 3, 4, 6\}$ replicates, and large-scale algorithmic and tissue-replicate studies report on the order of tens of replicate-like inputs.\footnote{Replicate counts are approximate and vary by experiment within each consortium; we use indicative ranges drawn from consortium release notes.} Aggregating bivariate IDR pairwise across replicate pairs can distort calibration or inflate downstream FDR unless the aggregation rule is itself modeled. Second, the reproducible component imposes an exchangeable bivariate Gaussian dependence with a single correlation parameter; consortium data with heterogeneous platforms may exhibit factor-like correlation patterns that are not exchangeable, and forcing exchangeability can bias local-idr estimates on platform-discordant features. Third, the bivariate Gaussian latent structure imposes a single shared marginal Gaussian latent variance across replicates after pseudo-data transformation; replicates from different sequencing depths, platforms, or tissue types can have different marginal score distributions, so a shared-variance assumption may miscalibrate tail-region idr estimates exactly where reproducibility decisions are made.

\begin{takeawaybox}
The original IDR was designed for two replicates with shared latent dependence and shared marginals. PY-IDR relaxes the corresponding limitations through three matched design choices: a Pitman--Yor mixture for multiple reproducible dependence regimes, structured per-cluster correlation matrices for multi-replicate dependence, and replicate-specific Bernstein--Dirichlet priors for heterogeneous marginals.
\end{takeawaybox}

\subsection{Why Pitman--Yor}

A natural Bayesian extension uses a nonparametric mixture over copula structures. The Pitman--Yor process \citep{PitmanYor1997} generalizes the Dirichlet process by introducing a discount parameter $\sigma \in [0,1)$, yielding power-law cluster-size distributions that match heavy-tailed empirical regimes. The operative property motivating PY for our application is the growth rate of the expected number of distinct components: under a Dirichlet process with concentration $\alpha$ and $n$ observations, the expected number of occupied clusters grows on the order of $\alpha \log n$, while under a Pitman--Yor process with discount $\sigma > 0$ it grows on the order of $n^\sigma$ up to constants depending on $(\alpha,\sigma)$ \citep{PitmanYor1997}. Thus, at consortium-scale sample sizes, the prior can allocate non-negligible mass to many rare dependence regimes rather than forcing a small fixed set of clusters.

We extend the PY mixture framework from densities on Euclidean space to mixtures of copulas with structured correlation and planned tail-dependence atoms. PY mixtures of densities are well-studied with $L^1$-consistency results \citep{IshwaranJames2001,WuGhosal2010} and adaptive contraction theory in $L^p$-metrics \citep{Scricciolo2014BAPY}; the corresponding consistency theory for PY mixtures of \emph{copulas} requires additional boundary and identifiability conditions. The closest published competitor is the work of \citet{PanNietoBarajasCraiu2024}, who develop a Bayesian nonparametric mixture of Archimedean copulas with a Pitman--Yor mixing distribution. Our contribution differs from theirs in three central respects: we handle Gaussian copulas with structured correlation matrices rather than Archimedean families with scalar dependence parameters, we focus on the multi-replicate IDR application rather than dependence modeling per se, and we make explicit which contraction and mFDR-control statements are established under stated assumptions and which pieces remain proof obligations for the fully general PY-copula setting.

\subsection{Contributions}

This paper makes the following contributions, which we summarize briefly here and develop in subsequent sections.
\begin{itemize}[leftmargin=*]
\item We propose a Pitman--Yor copula-mixture model for the reproducible component of an IDR mixture, accommodating arbitrary $K \geq 2$, structured correlation across mixture atoms (exchangeable, one-factor, two-factor), and planned heavy-tail atomic types ($t_5$, Clayton, Gumbel) within a single nonparametric prior.
\item We add replicate-specific Bernstein--Dirichlet marginal priors that decouple the marginal models across replicates while retaining a tractable route to marginal approximation and contraction analysis.
\item We provide two complementary inference algorithms: a Pitman--Yor P\'olya-urn MCMC scheme using Neal-style auxiliary atoms \citep{Neal2000} and the predictive structure of \citet{Pitman1996BlackwellMacQueen,IshwaranJames2001}, and a structured mean-field variational algorithm with stochastic natural-gradient updates in the framework of \citet{HoffmanBleiWangPaisley2013}.
\item We give a theoretical analysis of finite-sieve identifiability, conditional posterior contraction, asymptotic Bayes-mFDR control, Harris ergodicity, and the additional assumptions needed for stronger PY-tail, geometric-ergodicity, and variational-rate claims.
\item We design a planned empirical protocol benchmarking PY-IDR against vanilla IDR \citep{LiBrownHuangBickel2011}, eCV \citep{Iakovidis2024eCV}, nestedIDR \citep{Ranalli2026nestedIDR}, MaRR \citep{Philtron2018MaRR}, IDR2D \citep{Krismer2020IDR2D}, ChIP-R \citep{Newell2021ChIPR}, and the BNP Archimedean copula mixture of \citet{PanNietoBarajasCraiu2024} on five real datasets and an extended S1--S5 simulation suite.
\end{itemize}

\paragraph{Design map.}
The three modeling ingredients map directly to the three limitations above. Multi-replicate dependence is handled by $K$-dimensional copulas and structured correlation classes. Heterogeneous marginals are handled by replicate-specific Bernstein--Dirichlet priors rather than a shared latent Gaussian marginal. Unknown and rare dependence regimes are handled by the Pitman--Yor mixing distribution, whose discount parameter permits heavier cluster-size tails than the Dirichlet-process special case.

\paragraph{Relation to prior work.}
Among prior multi-replicate IDR variants, our method is intended to combine arbitrary replicate count, latent dependence modeling, heterogeneous non-Gaussian marginals, posterior uncertainty, and explicit theoretical analysis in a single Bayesian framework. We avoid relying on unverified novelty claims here; a final submission should verify novelty and implementation availability against the most recent software and bibliographic sources. Table~\ref{tab:related-work} in the next section summarizes the principal comparators along the dimensions of replicate count, copula family, marginal model, FDR-control mechanism, and theoretical guarantees.
